% =============================================================================
% 5.3 Documentacion y Cobertura de la API REST
% Archivo: 5.3_cobertura_api.tex
% =============================================================================

\section{Documentaci\'on y Cobertura de la API REST}
\label{sec:cobertura-api}

La documentaci\'on exhaustiva de la interfaz de programaci\'on de aplicaciones
(\textit{Application Programming Interface}, API) constituye un pilar
fundamental en el desarrollo de sistemas distribuidos, puesto que establece el
contrato formal entre el servidor \textit{backend} y los clientes consumidores.
La plataforma implementa una API REST documentada bajo la especificaci\'on
OpenAPI 3.0, la cual provee una descripci\'on estructurada y legible por
m\'aquinas de todos los \textit{endpoints} disponibles, sus par\'ametros de
entrada, esquemas de respuesta y mecanismos de autenticaci\'on.

\subsection{Especificaci\'on T\'ecnica de la API}
\label{subsec:especificacion-api}

El servidor de la API REST opera sobre el \textit{framework} Express.js,
expuesto en el puerto \texttt{8787} del entorno local de desarrollo
(\texttt{localhost}). La autenticaci\'on de las solicitudes se realiza mediante
tokens JWT (\textit{JSON Web Tokens}) transmitidos a trav\'es del esquema
\texttt{Bearer} en la cabecera \texttt{Authorization} de cada petici\'on HTTP.
La Tabla~\ref{tab:especificacion-api} sintetiza los par\'ametros t\'ecnicos
principales de la configuraci\'on de la API.

\begin{table}[htbp]
\centering
\caption{Especificaci\'on t\'ecnica de la API REST}
\label{tab:especificacion-api}
\begin{tabular}{|l|l|}
\hline
\textbf{Par\'ametro} & \textbf{Valor} \\
\hline
Especificaci\'on          & OpenAPI 3.0 \\
\hline
\textit{Framework}        & Express.js (Node.js) \\
\hline
URL base                  & \texttt{http://localhost:8787} \\
\hline
Esquema de autenticaci\'on & Bearer JWT \\
\hline
Formato de intercambio    & JSON (\texttt{application/json}) \\
\hline
Versionamiento            & Por prefijo de ruta (\texttt{/api/v1/}) \\
\hline
\end{tabular}
\end{table}

\subsection{Distribuci\'on de Endpoints por M\'odulo}
\label{subsec:endpoints-modulo}

La API se organiza en m\'odulos funcionales cohesivos que reflejan la
arquitectura de dominios del sistema. Cada m\'odulo agrupa un conjunto de
\textit{endpoints} relacionados sem\'anticamente, lo cual facilita el
mantenimiento, la documentaci\'on y la asignaci\'on de permisos granulares.
La Tabla~\ref{tab:endpoints-modulo} presenta la distribuci\'on de
\textit{endpoints} por m\'odulo funcional.

\begin{table}[htbp]
\centering
\caption{Distribuci\'on de \textit{endpoints} por m\'odulo de la API REST}
\label{tab:endpoints-modulo}
\begin{tabular}{|l|c|p{7cm}|}
\hline
\textbf{M\'odulo} & \textbf{Endpoints} & \textbf{Descripci\'on funcional} \\
\hline
\texttt{auth}       & var.\ & Autenticaci\'on, emisi\'on y renovaci\'on de tokens JWT, verificaci\'on OTP y evaluaci\'on de riesgo de acceso. \\
\hline
\texttt{audit}      & 2  & Registro y consulta de eventos de auditor\'ia del sistema, incluyendo accesos a datos sensibles. \\
\hline
\texttt{iam}        & 10 & Gesti\'on de identidad y acceso: roles, permisos, asignaci\'on de usuarios, y supervisi\'on de sesiones activas. \\
\hline
\texttt{onboarding} & 2  & Registro de organizaciones (\textit{tenants}), validaci\'on de RUC y creaci\'on de cuentas administradoras. \\
\hline
\texttt{sedes}      & 4  & Administraci\'on de sedes f\'isicas: registro, edici\'on, desactivaci\'on y consulta de locales operativos. \\
\hline
\hline
\textbf{Total}      & \textbf{18+} & Cobertura completa de los dominios funcionales cr\'iticos del sistema. \\
\hline
\end{tabular}
\end{table}

\subsection{An\'alisis de la Documentaci\'on OpenAPI}
\label{subsec:analisis-openapi}

La adopci\'on de la especificaci\'on OpenAPI 3.0 como est\'andar de
documentaci\'on confiere m\'ultiples ventajas al proceso de desarrollo. En
primer lugar, permite la generaci\'on autom\'atica de interfaces interactivas
de exploraci\'on mediante herramientas como Swagger UI, lo cual acelera la
integraci\'on de los equipos de \textit{frontend} y reduce la dependencia de
documentaci\'on manual propensa a desactualizaci\'on. En segundo lugar,
habilita la validaci\'on autom\'atica de las solicitudes y respuestas contra
los esquemas definidos, proporcionando una capa adicional de aseguramiento de
calidad en tiempo de ejecuci\'on.

\begin{figure}[htbp]
\centering
\includegraphics[width=0.95\textwidth]{coberapi}
\caption{Interfaz de documentaci\'on interactiva de la API REST generada mediante OpenAPI 3.0}
\label{fig:cobertura-api}
\end{figure}

La Figura~\ref{fig:cobertura-api} muestra la interfaz de documentaci\'on
interactiva generada a partir de la especificaci\'on OpenAPI. Se observa la
organizaci\'on jer\'arquica de los \textit{endpoints} por m\'odulo funcional,
los esquemas de solicitud y respuesta, y los c\'odigos de estado HTTP asociados
a cada operaci\'on. Esta documentaci\'on se genera de forma din\'amica a partir
del c\'odigo fuente del servidor, garantizando la consistencia entre la
implementaci\'on y la especificaci\'on documentada.

El m\'odulo \texttt{iam}, con diez \textit{endpoints}, constituye el componente
de mayor complejidad funcional, dado que encapsula la totalidad de las
operaciones de gesti\'on de identidad y control de acceso basado en roles
(RBAC). Este m\'odulo implementa operaciones CRUD sobre roles y permisos,
as\'i como la asignaci\'on contextual de usuarios a sedes f\'isicas
espec\'ificas, lo cual refleja la naturaleza multi-sede de las organizaciones
no gubernamentales atendidas por la plataforma.
